\scanpage{11}
\deleted{notant que}
en particulier, tenant compte de $3^\circ$,
\diagram{\begin{tikzcd}[ampersand replacement=\&]
\mathrm{I}^{\unsure{a|\alpha}}(G) \arrow[r] \& \mathrm{Pol}(G)\times\mathrm{Eff}(G)
\end{tikzcd}}
\inserted{les $i:G_{m\mathbb{C}}\longrightarrow G_{\mathbb{C}}$ définissant une polarisation, modulo action de $G(\mathbb{R})$\ldots}

\uline{Questions} Cette dernière application, \inserted{$\mathrm{I}(G)-\mathrm{Eff}(G)$, injective ?}

Les structures $\in\mathrm{Eff}(G)$ provenant de la théorie des motifs peuvent elles \unsure{tjs} se décrire pr un $i_1$ ? i.e. sont elles dans $\operatorname{Im}\mathrm{I}(G)$ ?

\margin{N.B. Le fait qu'une structure d'effectivité puisse se décrire pr un élément de $\mathrm{I}(G)$ (resp $\mathrm{I}^{\unsure{a|\alpha}}(G)$) est une propriété intrinsèque à la catégorie de motifs $M$, indép. du foncteur qui la réalise. ! !}

(Regarder ce qui s'obtient en tordant pr un torseur $\xi$. Examiner aussi, pr ces questions et celles posées dans $3^\circ$, le cas où $G$ avec ses $4$ structures provient d'un motif sur schéma \inserted{connexe} $S$ muni d'un $\xi\in S(\mathbb{Q})$, en utilisant le foncteur de De-Rham\ldots\ ?]

Regarder $G$ commutatif ! $G=\unsure{\mathrm{Sl}(2,\mathbb{Q})|\mathrm{Gl}(2,\mathbb{Q})}$ !

Des propriétés arithmétiques plus profondes peuvent être sans doute formulées pr l'introduction des éléments de Frobenius et l'interprétation des structures $3^\circ$ et $4^\circ$ en termes de ces derniers.

\mspage{10}
